\documentclass[11pt,a4paper]{article}
\usepackage[english]{babel}
\usepackage{graphicx}
\usepackage{amsmath}
\usepackage{amssymb}
\usepackage{hyperref}
\usepackage{geometry}
\usepackage{float}
\usepackage{listings}
\usepackage{xcolor}
\usepackage{enumitem}
\usepackage{booktabs}
\usepackage{setspace}
\onehalfspacing
\usepackage{parskip}
\setlength{\marginparwidth}{2cm}
\usepackage{todonotes}
%\usepackage[disable]{todonotes}
\usepackage{csquotes}
\usepackage[style=numeric, backend=biber]{biblatex}
\addbibresource{refs.bib}
\usepackage[export]{adjustbox}
\usepackage{xcolor}
\usepackage{subcaption}

\usepackage{titlesec}
\newcommand{\sectionbreak}{\clearpage}

% find all images in this folder
\graphicspath{ {./img/} }

\geometry{margin=1in}

\lstset{
    basicstyle=\ttfamily\small,
    breaklines=true,
    frame=single,
    language=Python
}

\title{\textbf{OpTime: Optimising Timetables} \\ \large Project Report: UI/UX, Gurobi, and Simulated Annealing}

\author{
    \textbf{Front-End \& UI/UX Group} \\
    Haonan Wu (126815), Ruyue Zhao (127561) \\[1ex]
    \textbf{Gurobi Optimization Group} \\
    Polina Evgrafova, Catalin Viscun (125396)\\[1ex]
    \textbf{Simulated Annealing Group} \\
    Abenezer Belay (128249), Chaithanya Thalapaneni  (127081)
}
\date{WiSe 2025 \\ Advisor: Andreas Jakoby}

\begin{document}

\maketitle
\newpage
\tableofcontents
\newpage

\section{Introduction}
Scheduling university timetables is a fundamental yet demanding administrative task. It requires balancing the competing needs of various stakeholders, including students, lecturers, and departmental offices. In many faculties, this process is still handled manually by administrative staff---a task that is both mentally exhausting and complex. Creating a feasible schedule that satisfies everyone while respecting resource limits is a significant challenge. Our project aims to alleviate this administrative burden by integrating optimization algorithms into a user-friendly, web-based system.

Timetable scheduling is a computationally hard problem. The number of possible combinations between courses, rooms, time slots, and lecturers grows exponentially as data increases. Attempting to find a valid solution manually is not only slow but also highly error-prone.

To address this, we developed a web-based tool specifically tailored for the Faculty of Media. The project was carried out by three sub-teams focusing on different core components:

\begin{itemize}
    \item \textbf{Frontend (UI/UX)}: Built the user interface for managing lecturers, courses, and rooms, and for visualizing the generated timetable.
    \item \textbf{Gurobi}: Used Integer Linear Programming to generate an initial valid schedule that satisfies all hard constraints.
     \item \textbf{Simulated Annealing}: Applied a metaheuristic algorithm to improve the current schedule by reducing soft constraint violations.
\end{itemize}

The three parts work together in a user-driven process: the frontend collects all input data and sends it to the Gurobi solver to create a base schedule. Users can then manually adjust the results through the interface. After that, the Simulated Annealing module takes that version and improves it further. The final result is shown on the frontend, ready to be reviewed and exported. This document covers each component in detail.

\section{Literature Review}
Before we started building the OpTime System, our group analyzed 13 research papers to study the principles of various optimization algorithms and their practical applications in university timetabling. This research gave us a clear picture of the technical challenges and helped us decide which algorithms to use for our specific scheduling needs.

\subsection{Problem Complexity}
The literature defines university timetabling as an NP-complete or NP-hard problem, meaning no efficient algorithm can find a perfect solution every time as the data size grows \cite{burke1997automated, schaerf1999survey}. To address this issue, exhaustive enumeration of all feasible configurations in search of an optimal one is avoided. Instead, a suboptimal solution is sought. Several approaches exist for this purpose, including, among others\cite{alghamdi2020review}:

\begin{itemize}
    \item The use of off-the-shelf solvers (engines) based on integer programming, where the problem is formulated by specifying constraints interpretable by the model in such a way that the solution to the formulated problem corresponds to a feasible solution of the original problem.
    
    \item Algorithms such as simulated annealing, implemented either using existing frameworks or developed from scratch, which iteratively select improved configurations from a set of modifiable candidates, exclude certain configurations according to problem-specific constraints, and evaluate admissible configurations using an objective function.

    \item The use of genetic algorithms, which mimic key aspects of the evolutionary process. In this approach, new configurations are generated from existing ones, making it conceptually related to the method described above, while also incorporating mechanisms for repairing infeasible solutions. By imposing limits on the number of iterations or steps in these and the previously mentioned algorithms, it is possible to obtain a suboptimal yet satisfactory solution within a reasonable time frame.\cite{Kristiadi2019}
    
    \item Hybrid approaches that combine multiple algorithms, distributing tasks between them in order to avoid evaluating all possible configurations, while compensating for this by enabling transitions from inferior to superior configurations.\cite{Yunita2018}
\end{itemize}

\subsection{Algorithm Mechanisims and Applications}
Our review focused on the mechanics of several optimization methods:
\begin{itemize}
    \item \textbf{Genetic Algorithms (GA)}: These use a biological evolution approach involving selection, crossover, and mutation. We found that Violated Directed Mutation (VDM) is particularly effective because it specifically targets and fixes genes that violate constraints rather than making random changes \cite{Kristiadi2019}.
    \item \textbf{Simulated Annealing (SA)}: This algorithm keeps track of a current solution and occasionally accepts ``worse'' moves as the system ``temperature'' decreases. This mechanism helps the solver escape local optima to find a better overall schedule. \cite{jakoby2026discrete}
    \item \textbf{Ant Colony Optimization (ACO) \& PSO}: ACO mimics ant behavior by using pheromones to find the best paths to a solution. When hybridized with Particle Swarm Optimization (PSO), the system can automatically tune its own parameters, making the search process more stable and successful \cite{Yunita2018}.
    \item \textbf{Fuzzy Logic}: Some models use fuzzy logic to handle uncertain human preferences, which helps in balancing conflicting soft constraints \cite{robinson2025adaptive}.
\end{itemize}

These insights led us to adopt a dual-solver architecture. While Integer Programming (IP) models can provide near-optimal solutions, in practice such solutions may fail to meet all practical requirements. This motivates the need for a mechanism to modify the obtained solution while preserving its feasibility and overall quality. The first component of the approach, due to its deterministic nature and lack of stochasticity, is not capable of providing such flexibility, thereby necessitating a complementary method.

We addressed these performance constraints by adopting a dual-solver architecture. We use Simulated Annealing for the drafting and real-time editing phases, as it provides the immediate feedback necessary for an interactive tool. This hybrid setup allows OpTime to stay fast and responsive at scale, while ensuring all results remain grounded in the formal logic of the IP-based model \cite{strichman2017near}.

% ==========================================
% PART I: Front-end, UI/UX
% ==========================================
\section{Part I: Front-End Design \& Implementation}
The front end serves as the central interface of the OpTime system, acting as the bridge between administrators and the two backend solvers. All user inputs, including lecturer availability constraints, course configurations and room capacities etc., are entered here and translated into structured data for the optimization engines. Conversely, solver outputs from both Gurobi and Simulated Annealing are parsed and rendered as readable and interactive timetables.
\subsection{UI/UX Design \& User Experience}
\subsubsection{Information Hierarchy and Navigation}
OpTime prioritizes administrative efficiency by using a dashboard-centric interface. This layout minimizes task friction, ensuring that core navigation and data tools remain accessible from any view (Figure \ref{fig:dashboard}).
\begin{itemize}
    \item \textbf{Global Navigation}: The left-hand sidebar acts as a unified navigation hub, providing direct access to Lecturers, Courses, and Rooms. This design maintains user orientation across the platform's data hierarchy.
    \item \textbf{Version Control}: The ``Saved Schedules'' repository allows administrators to archive and compare previous generated timetables. Each entry is auto-stamped to allow users to test different constraints and roll back to an ``optimal'' configuration via the Restore Data function.
    \item \textbf{Data Grid System}: The system employs high-density tables for efficient data management. Instead of permanent deletion, a status toggle (``Active'' vs. ``Inactive'') allows administrators to temporarily exclude specific entries from the solver's scope. This ``soft disable'' approach ensures that historical data remains intact for future use without cluttering the current optimization process.
\end{itemize}

\begin{figure}[H]
    \centering
    \includegraphics[width=0.8\textwidth, cfbox=lightgray 0.5pt 0pt]{Documentation/img/ui_main_dashboard.png}
    \caption{OpTime dashboard: schedule versioning and data management.}
    \label{fig:dashboard}
\end{figure}

\subsubsection{The Interactive Availability Matrix}
To bridge the gap between faculty preferences and the automated scheduler, we developed an intuitive Availability Matrix (Figure \ref{fig:matrix}).
Instead of asking users to fill out complex forms, we designed a 7-day visual grid divided into 90-minute blocks. The interface covers the core academic window (09:15--18:30) while also providing extended slots at 07:30--09:00 and 19:00--20:30. These additions allow the system to handle early or late sessions required by specific departments.

The system uses a straightforward four-color code to define how the optimization engine should handle different time slots:

\begin{itemize}
\item \textbf{Red (Unavailable)}: Acts as a hard boundary. The solver is strictly prohibited from scheduling any course or lecturer in these time slots.
\item \textbf{Blue (Preferred)}: Represents the preferred time slots for the lecturer or the specific course.
\item \textbf{Green (Available)}: Serves as a standard fallback. If the preferred blue slots cannot be met due to scheduling conflicts, the system uses these areas to ensure the final schedule remains valid.
\item \textbf{Yellow (Rather Not)}: Identifies a negative soft preference. The algorithm tries to avoid these slots whenever possible, utilizing them only as a last resort to find a feasible solution.
\end{itemize}

To minimize redundant data entry, the system includes a Smart Syncing mechanism. Manually inputting the same availability for multiple courses is inefficient and often leads to errors. We addressed this by synchronizing the data flow between lecturers and their assigned courses.

Through the ``Copy to Assigned Courses'' feature, a lecturer’s general availability can be pushed to all their assigned classes in one action. Conversely, when configuring a specific course, the system can automatically pull the assigned lecturer's schedule. This ensures data consistency across the platform and prevents scheduling conflicts, such as placing a class during a time when the lecturer is marked as unavailable.

\begin{figure}
    \centering
    \begin{subfigure}[b]{0.8\textwidth}
        \centering
        \includegraphics[width=\textwidth, cfbox=lightgray 0.5pt 0pt]{Documentation/img/ui_availability_lecturer1.jpeg}
        \caption{Lecturer's view of availability and preferences.}
        \label{fig:matrix_l}
    \end{subfigure}

    \vspace{2em}

    \begin{subfigure}[b]{0.8\textwidth}
        \centering
        \includegraphics[width=\textwidth, cfbox=lightgray 0.5pt 0pt]{Documentation/img/ui_availability_course1.jpg}
        \caption{Course-specific time slot constraints.}
        \label{fig:matrix_c}
    \end{subfigure}
    \caption{Visual representation of time slot availability and preferences.}
    \label{fig:matrix}
\end{figure}

\subsubsection{Data Entry and Entity Management}
To maintain a continuous workflow, the system handles data entry for lecturers, courses, and rooms through modal windows (Figure \ref{fig:course_modal}). This design allows users to add or edit information without navigating away from the main list, reducing unnecessary page reloads.
\begin{itemize}
    \item \textbf{Entity-Specific Data Entry}: The system provides specific modals for managing lecturers, courses, and rooms. While entries for lecturers and rooms focus on basic attriubtes such as faculty and capacity, the ``Add Course'' interface handles more complex configurations. It allows users to link multiple lecturers to a single course and define curriculum requirements (e.g., mandatory vs. elective). This design ensures that the front-end captures all the parameters required by the Gurobi and SA solvers.
    \item \textbf{Input Validation and Integrity}: To prevent invalide data from reaching the optimization engine, the UI performs real-time validation. It checks for formatting errors and missing fields before any data is committed to the database, ensuring that the solver always oprates on a consistent and high-quality dtaset.
\end{itemize}

\begin{figure}
    \centering
    \begin{subfigure}[b]{0.8\textwidth}
        \centering
        \includegraphics[width=\textwidth, cfbox=lightgray 0.5pt 0pt]{Documentation/img/ui_add_course1.png}
        \caption{Initial view of the modal.}
        \label{fig:course_modal1}
    \end{subfigure}

    \vspace{2em}

    \begin{subfigure}[b]{0.8\textwidth}
        \centering
        \includegraphics[width=\textwidth, cfbox=lightgray 0.5pt 0pt]{Documentation/img/ui_add_course2.png}
        \caption{Scrolled view showing remaining fields.}
        \label{fig:course_modal2}
    \end{subfigure}
    \caption{The ``Add Cours'' modal interface (scrolled to display all fields).}
    \label{fig:course_modal}
\end{figure}

\subsubsection{Visualization and Interaction of Generated Timetables}
Once the optimization is complete, the system converts raw solver data into an intuitive weekly planner (Figure \ref{fig:timetable}) with the following features:
\begin{itemize}
    \item \textbf{Layout and Information Display}: The generated timetable uses a standard weekly grid, as shown in Figure\ref{fig:timetable_view}. Each session is rendered as a color-coded card displaying essential metadata: the course name, lecturer, and room code. This allows users to grasp the overall schedule at a glance.
    \item \textbf{Manual Adjustments and Visual Feedback}: The system supports a drag-and-drop interface for manual fine-tuning (Figure\ref{fig:timetable_adjust}). To prevent users from inadvertently breaking the scheduling logic, the UI provides real-time visual cues: valid time slots are highlighted in green, while prohibited slots (due to lecturer or room conflicts) are marked in red. This feedback loop guides users to make adjustments that remain within the bounds of the system's constraints.
    \item \textbf{Conflict Reporting}: A ``View Report'' button allows users to access a detailed diagnostic modal (Figure\ref{fig:timetable_report}). Unlike existing timetabling interfaces that leave constraint interpretation entirely to the user \cite{mccollum2014review}, the report categorizes findings into ``Blocking Issues'' (hard constraint violations) and ``Suggestions'' (soft constraint optimizations, such as inefficient room usage). This report helps users review the solver's logic and identify exactly where manual adjustments are needed.
\end{itemize}

\begin{figure}
    \centering
    \begin{subfigure}[b]{0.8\textwidth}
        \centering
        \includegraphics[width=\textwidth, cfbox=lightgray 0.5pt 0pt]{Documentation/img/ui_generated_timetable.png}
        \caption{Generated weekly timetable.}
        \label{fig:timetable_view}
    \end{subfigure}

    \vspace{1em}

    \begin{subfigure}[b]{0.8\textwidth}
        \centering
        \includegraphics[width=\textwidth, cfbox=lightgray 0.5pt 0pt]{Documentation/img/ui_timetable_adjust.png}
        \caption{Drag-and-drop with conflict highlighting.}
        \label{fig:timetable_adjust}
    \end{subfigure}

    \vspace{1em}

    \begin{subfigure}[b]{0.8\textwidth}
        \centering
        \includegraphics[width=\textwidth, cfbox=lightgray 0.5pt 0pt]{Documentation/img/ui_timetable_report.jpeg}
        \caption{Constraint violation report.}
        \label{fig:timetable_report}
    \end{subfigure}
    
    \caption{Timetable viewing, adjustment, and reporting workflow.}
    \label{fig:timetable}
\end{figure}

\subsection{Technical Implementation}

\subsubsection{Front-End Architecture}

The front-end system is implemented as a \textbf{modular single-page application (SPA)} that integrates multiple functional views within a unified interface. The architecture is designed to separate concerns between data management, interaction logic, and visualization, ensuring both scalability and maintainability.

\begin{itemize}
    \item \textbf{Modular Structure}: The system is decomposed into independent feature modules (e.g., CourseModule, LecturerModule, RoomModule, TimetableModule), each responsible for a specific domain. These modules are initialized and coordinated through a central entry point, enabling loose coupling and extensibility.
    
    \item \textbf{Centralized Data Management}: A singleton \textbf{DataManager} maintains a consistent global state across all modules. This ensures that updates in one component (e.g., editing a course or lecturer) are immediately reflected throughout the system.
    
    \item \textbf{Dynamic View Switching}: The interface organizes content into multiple views (Lecturers, Courses, Rooms, Timetable), which are dynamically activated without page reload. This allows seamless navigation while preserving user context and interaction state.
    
    \item \textbf{Reusable UI Components}: HTML partials are loaded dynamically at runtime, enabling flexible composition of interface elements such as modals, drawers, and tool panels. This reduces redundancy and simplifies future interface extensions.
    
    \item \textbf{Structured Layout Design}: The interface combines a persistent sidebar for global actions (navigation, generation, data management) with a main content area for task-specific interaction. This layout supports efficient workflows and minimizes user disorientation.
    
\end{itemize}
This modular architecture improves maintainability and scalability, while ensuring a clear separation of concerns.
It also supports a responsive and predictable interface, which enhances usability and user experience.



\subsubsection{The Dynamic Timetable Engine}

The core functionality of the system is implemented through a \textbf{dynamic timetable engine} that integrates algorithmic scheduling results with interactive front-end manipulation. This component transforms backend optimization outputs into a structured, user-controllable visual representation.

\begin{itemize}
    \item \textbf{State-Driven Rendering}: The timetable is generated from structured assignment data returned by the backend solver. Each entry is mapped to a visual ``Course Card'', whose position is determined by its assigned day and time-slot index. The system maintains an internal state representation of all assignments, ensuring consistency between data and visualization.
    
    \item \textbf{Interactive Visualization Layer}: The timetable is rendered as a grid aligned with temporal structure (days and time slots), allowing users to intuitively interpret scheduling outcomes. Each grid cell dynamically aggregates multiple assignments, supporting real-time updates and filtering (e.g., by faculty or semester).
    
    \item \textbf{Real-Time Conflict Visualization}: The system integrates a validation layer that highlights constraint violations directly within the interface. Critical conflicts (e.g., lecturer overlap, room capacity issues) are encoded using red visual indicators, while softer constraint violations are represented using warning colors. This enables immediate identification of problematic assignments.
    
    \item \textbf{Drag-and-Drop Interaction}: Users can manually adjust scheduling results by dragging course cards across the grid. The system provides visual hints about slot suitability (available, preferred, or unavailable) and prevents invalid placements. Upon each drop action, the system updates the internal state and triggers validation to maintain constraint consistency.
    
    \item \textbf{Iterative Update Mechanism}: The engine supports continuous refinement by allowing users to regenerate or improve schedules based on current assignments. User modifications are preserved and re-integrated into subsequent optimization steps, enabling an interactive exploration of alternative scheduling solutions.
\end{itemize}

To ensure a responsive and uninterrupted user experience, the system employs \textbf{asynchronous operation handling} for data loading, scheduling computation, and system interaction.

\begin{itemize}
    \item \textbf{Asynchronous Initialization}: During application startup, HTML partials and initial data are loaded using asynchronous operations. The system waits for all components to be successfully initialized before rendering the interface, ensuring consistency while avoiding blocking the main thread.
    
    \item \textbf{Promise-Based Backend Communication}: The front-end communicates with the scheduling solver through \texttt{fetch}-based API calls. This allows timetable generation and improvement processes to run in the background while the user interface remains interactive.
    
    \item \textbf{Loading State Feedback}: During computationally intensive operations (e.g., schedule generation or improvement), the system displays a loading overlay with status messages (e.g., ``Optimizing schedule'' or ``Improving schedule''). This provides clear feedback about system status and reduces user uncertainty.
    
    \item \textbf{Non-Blocking Interaction Model}: User interactions such as editing assignments, dragging course cards, or switching views are handled without interrupting the application flow. Updates are applied immediately to the local state, while validation and backend processing occur asynchronously.
    
    \item \textbf{Selective Rendering}: The system conditionally renders content based on the active view, avoiding unnecessary updates and improving performance. This ensures that only relevant components are updated in response to user actions.
\end{itemize}

\subsection{Data Interaction \& Integration}

\subsubsection{Bridge Between User and Algorithm}

The system establishes a \textbf{bidirectional interaction bridge} between user input and the scheduling algorithm, transforming abstract optimization constraints into an interactive and controllable interface.

\begin{itemize}
    \item \textbf{Explicit Constraint Specification}: Users define scheduling constraints through an interactive grid interface, where each time slot can be assigned one of three states: available, preferred (soft constraint), or unavailable (hard constraint). These selections are directly translated into structured constraint data used by the optimization engine.
    
    \item \textbf{Unified Data Representation}: All user-defined inputs (courses, lecturers, rooms, and time constraints) are collected and managed through a centralized data layer. This ensures consistency between user actions and algorithmic computation.
    
    \item \textbf{Bidirectional Interaction}: Users can both initiate automated scheduling (via ``Generate'' or ``Improve'') and manually adjust results (e.g., drag-and-drop or edit dialogs). Each interaction updates the system state and can be re-evaluated by the backend solver.
    
    \item \textbf{Human-in-the-Loop Optimization}: The system supports an iterative cycle in which users refine schedules based on algorithmic feedback. User modifications are preserved and re-submitted to the solver, enabling continuous improvement of scheduling outcomes.
\end{itemize}


\subsubsection{Intelligent Conflict Reporting}

To support user understanding and decision-making, the system provides \textbf{multi-level conflict detection and feedback} mechanisms.

\begin{itemize}
    \item \textbf{Constraint Classification}: Detected violations are categorized into two types:
        \begin{itemize}
            \item \textbf{Critical Issues (Hard Constraints)}: Logical conflicts that must be resolved (e.g., lecturer overlap, room capacity violations).
            \item \textbf{Preference Warnings (Soft Constraints)}: Non-critical deviations from preferred conditions (e.g., suboptimal time slots or room assignments).
        \end{itemize}
    
    \item \textbf{Visual Encoding}: Conflicts are directly visualized within the timetable using color-coded indicators (e.g., red for critical conflicts and orange for warnings). This allows users to quickly identify problematic assignments without inspecting raw data.
    
    \item \textbf{Detailed Reporting}: A dedicated report interface presents structured explanations of conflicts, including their causes and suggested corrective actions. This enhances interpretability and supports efficient problem resolution.
    
    \item \textbf{Real-Time Validation}: After each user interaction (e.g., manual edits or drag-and-drop operations), the system triggers validation processes to ensure that the current schedule remains consistent with defined constraints.
\end{itemize}


\subsubsection{Optimization Workflow \& Persistence}

The system supports an \textbf{iterative and mixed-initiative workflow} that combines automated optimization with user-driven refinement.

\begin{itemize}
    \item \textbf{Iterative Scheduling Process}: Users first generate an initial timetable, then refine it through manual adjustments or by invoking the ``Improve'' function. This enables exploration of multiple feasible solutions.
    
    \item \textbf{Interactive Manipulation}: Users can directly modify scheduling results through drag-and-drop and editing interfaces. These actions immediately update the system state and trigger re-validation.
    
    \item \textbf{Version Management}: The system provides functionality for saving, loading, and managing different timetable versions. This supports comparison between alternative solutions and prevents loss of progress.
    
    \item \textbf{Data Export and Persistence}: Scheduling results can be exported (e.g., as PDF) and stored persistently. Backup and restore mechanisms ensure data safety and support long-term use.
\end{itemize}

Overall, the system transforms a complex optimization problem into an interactive, transparent, and user-controllable process, improving usability, interpretability, and user trust in algorithmic decision-making.

% ==========================================
% PART II: GUROBI
% ==========================================

\section{Part II: Optimization via Gurobi}
While the Simulated Annealing module provides a fast, stochastic approach to timetable optimization, the Gurobi module takes a complementary strategy: it formulates the scheduling problem as an Integer Linear Program (ILP) and solves it to provable optimality. Integer Linear Programming is a mathematical framework in which the goal is to minimize (or maximize) a linear objective function subject to a set of linear constraints, where some or all decision variables are restricted to integer values. For binary scheduling decisions---such as whether a course is assigned to a particular room, time slot, and lecturer---this framework is a natural fit\cite{Talmor2024OptimizingAT}.

Gurobi is a state-of-the-art commercial solver for mathematical optimization problems, widely used in both industry and academic research. Compared to open-source alternatives such as GLPK or CBC, Gurobi offers significantly faster solve times on large ILP instances, a clean Python API (gurobipy), and an academic license that makes it accessible for university projects. The key trade-off between the two approaches in OpTime is as follows: Gurobi guarantees finding the globally optimal schedule but it takes much longer as the problem size grows; Simulated Annealing handles large tasks much better but it doesn't guarantee a perfect result. Together, they give users a practical choice based on their specific needs.

The solution produced by Gurobi is obtained, as mentioned above, via linear integer programming (ILP), which is, in general, an NP-hard problem. Consequently, solving such models to optimality can be computationally challenging, especially for large-scale instances.

A widely used technique in this context is the relaxation of the model to a linear programming (LP) problem, in which the integrality constraints on the variables are temporarily removed, allowing them to take continuous values. Such LP problems can be solved in polynomial time using well-established algorithms.

The solution to the relaxed problem provides a bound on the optimal value of the original ILP and serves as a foundation for more advanced methods, such as branch-and-bound and branch-and-cut, which systematically explore the solution space while leveraging LP relaxations to efficiently prune infeasible or suboptimal regions.

%% написать что млп тоже нп хард и сказать что существуют оптимизации, без интеджеров легко - с интеджерами сложно, говорим что сначала считаем реалы потом идем к интеджерам
\subsection{Problem Formulation}
In order to solve the scheduling problem using the optimization software described above, it is necessary to construct a formal model that represents both the occurrence of a class session and the natural constraints imposed by the real-world environment in which such sessions take place. This representation must be expressed in a form that can be interpreted by a specific solver implementing a particular optimization algorithm.

The process of constructing this model can be conceptually divided into four stages, each corresponding to one of the fundamental components of the final optimization formulation. The first stage defines the decision variables, which form the foundation upon which all subsequent components of the model are built. The second stage introduces hard constraints, which establish the feasible boundaries of the problem and enforce the essential rules that must always be satisfied.The third stage incorporates soft constraints. Unlike hard constraints, these conditions represent preferences rather than strict requirements. They allow the model to be adjusted and adapted according to practical considerations and user needs.Finally, the fourth stage defines the objective function to be minimized. This function aggregates penalties associated with violations of soft constraints and provides the optimization solver with a quantitative criterion for evaluating different feasible schedules.

\subsubsection{Decision Variables}
The modeling process begins with the definition of decision variables as a multidimensional array of binary variables in Gurobi. Each unit entry in this array represents the assignment of a class to a specific time slot, taught by a particular instructor, for a given course, and in a designated room.

In Python-style notation, this can be expressed as:
\begin{verbatim}
x = m.addVars(slots, teachers_ids, courses_ids, rooms_ids, vtype=GRB.BINARY, name="x")
\end{verbatim}

Such an approach enables a natural and flexible representation of the system's constraints, as all relevant information is organized within data structures describing instructors, courses, and rooms. This facilitates the logical formulation of constraints over the entire space of possible assignments.
\subsubsection{Hard Constraints}
In the context of timetable generation, the core of the model is formed by \emph{hard constraints}, which define the set of feasible and infeasible configurations. In practical terms, these constraints correspond to physically impossible situations that must be excluded from the solution space.

Such constraints can generally be expressed in two principal ways. First, by enforcing that certain entities exist in at most one instance at any given time: 

\begin{verbatim}
for r in rooms:                                  
    for s in slots:                              
        m.addConstr(
            gp.quicksum(x[s, t, c, r]
                        for t in teachers
                        for c in courses) <= 1), ...)
\end{verbatim}

Second, constraints can enforce the complete absence of certain configurations:
\begin{verbatim}
for t_id in teachers_id_list:
    teacher = dict_teacher_id_obj[t_id]
    for slot in teacher.hard_time_constr:
        m.addConstr(
            gp.quicksum(
                x[slot, t_id, c_id, r_id]
                for c_id in courses_ids
                for r_id in rooms_ids
            ) == 0), ...)
\end{verbatim}

The meaning of these examples is as follows: at any given time, each room can host at most one class, taught by a single instructor for a single course; additionally, for any time slot excluded by a given instructor, there exists no room in which that instructor is assigned to teach any course.

As illustrated, once the problem is formulated in a form suitable for a solver, it becomes a powerful and flexible tool, allowing for a wide range of constraints to be incorporated and adjusted in accordance with real-world requirements.
\subsubsection{Soft Constraints}
Soft constraints are formulated in a similar manner; however, instead of explicitly forbidding certain configurations, additional variables and linear expressions are introduced to represent configurations that are undesirable to varying degrees.

For example, consider the following expression:
\begin{verbatim}
if room_by_id[r1].address != room_by_id[r2].address:
    expr += (
        x[s, t_id, c_id, r1]
        + x[s_next, t_id, c_id, r2]
    )

return weight * expr
\end{verbatim}

Such a linear expression becomes strictly positive whenever the desired condition is violated. It aggregates contributions from all configuration components in which the condition does not hold, thereby quantifying the extent of the violation.

In some cases, it is convenient to introduce auxiliary binary variables to explicitly model compound conditions:
\begin{verbatim}
if room_by_id[r_id].address in remote_buildings:
    z = m.addVar(vtype=GRB.BINARY)
    m.addConstr(z <= x[s1, t_id, c_id, r_id])
    m.addConstr(z <= x[s2, t_id, c_id, r_id])
    m.addConstr(z >= x[s1, t_id, c_id, r_id] + x[s2, t_id, c_id, r_id] - 1)
    expr += z

return weight * expr
\end{verbatim}

Here, the auxiliary variable $z$ captures the joint occurrence of two assignments, enabling a more precise formulation of the corresponding penalty term. The overall expression is then constructed as a weighted sum of such variables.

By formulating constraints of different types and relative importance in this manner, the model gains the ability to impose a hierarchy of preferences. This allows the solver not only to ensure feasibility with respect to hard constraints, but also to optimize solution quality by minimizing violations of soft constraints according to their assigned weights.
\subsubsection{Objective Function}
The aforementioned minimization constitutes the primary objective handled by Gurobi. In this setting, the solver is tasked, first, with satisfying all hard constraints, and second, with violating soft constraints only to the extent and in the manner that minimizes the total weighted penalty.

The assigned weights serve to express the relative importance of different constraints. For instance, low-priority preferences may include keeping classes in the same room; medium-priority constraints may enforce that instructors teach within their preferred time slots; and high-priority requirements may ensure that courses are held within the appropriate faculty and that mandatory and semi-mandatory courses do not overlap within a given program of study.

This weighted formulation enables a nuanced trade-off between competing objectives, guiding the solver toward solutions that are not only feasible but also aligned with practical and organizational priorities.

\subsection{Other implementation Details}
In addition to the modeling approach described above, several auxiliary components were implemented to improve usability and facilitate integration with the existing codebase. In particular, a central module was developed to define the interaction workflow between the constructed model, the solver, and the resulting data. This module is also responsible for collecting and exposing diagnostic information about the model, such as feasibility status (feasible/infeasible), execution time, and related metrics.

Furthermore, to enhance user interaction, an error and constraint-violation aggregator was implemented. This component provides structured feedback in cases where no feasible solution is found, enabling users to better understand the causes of infeasibility and to adjust input data or constraints accordingly.


\subsection{Gurobi installation}
The first step in the implementation process was obtaining an academic license for Gurobi via its official website:
\begin{verbatim}
https://www.gurobi.com/academia/academic-program-and-licenses/
\end{verbatim}
In our case, an individual student license was used. For deploying the solution in a production environment, a web license or a site license would likely be required.

Subsequently, Gurobi and its Python interface were installed using \texttt{pip}. The license setup and activation process is documented in the user account dashboard, under the ``Licenses'' section.
\subsection{Integration with the Frontend}
To provide an end-to-end workflow from data entry to an actual timetable, the project connects a browser-based frontend (JavaScript) with a FastAPI backend (Python) that runs the optimization model in Gurobi. This technology split was a small hurdle at first: the UI manages state and interactions in the browser, whereas the solver requires consistent, typed structures and stable integer identifiers. The integration resolves this by using a clear REST interface and an explicit transformation layer on the backend.

The frontend maintains courses, lecturers, and rooms in local arrays and synchronizes them with the backend via REST. On startup, the current dataset is loaded through \texttt{GET /api/data}; any modifications performed in the UI are persisted through \texttt{POST /api/data}. The backend stores the state in \texttt{data/database.json}, which keeps the development setup simple, transparent, and reproducible.

A central linking component is the weekly availability editor. The UI implements a fixed $7\times7$ grid using keys of the form \texttt{"dayIndex\_slotIndex"}, e.g.\ \texttt{"0\_0"} for Monday in the first slot or \texttt{"6\_6"} for Sunday in the last slot. In the UI, these selections are stored as lists of slot keys (\texttt{available\_slots} and optionally \texttt{unavailable\_slots}). In the backend, they are translated into the solver's constraint sets, most importantly \texttt{hard\_time\_constr} (forbidden slots) and, in principle, also \texttt{soft\_time\_constr} and preference sets if those model components are enabled. Concretely, if an \texttt{available\_slots} list is present, all non-selected cells are treated as unavailable (via complement construction) and are passed to the solver as hard time restrictions. This makes UI interactions in the grid directly enforceable as solver constraints.

Timetable generation is triggered via \texttt{POST /api/plan}. The backend pipeline reads the latest UI database snapshot, converts it into a canonical solver input format, and then runs Gurobi (optionally without an improvement stage). The conversion happens in the adapter \texttt{ui\_db\_to\_canonical\_input}, where (i) UI identifiers (e.g.\ \texttt{l\_5749}) are deterministically mapped to integer IDs, (ii) courses, lecturers, and rooms are converted into the internal data structures, and (iii) UI slot keys are translated into solver time indices.

For consistency, the backend generates 49 fixed time slots ($7$ days $\times$ $7$ slots) and provides bijective mappings between UI keys and \texttt{timeslot\_id}s:
\[
\texttt{"0\_0"} \mapsto 1,\ \ldots,\ \texttt{"6\_6"} \mapsto 49.
\]
This mapping ensures that time information coming from the UI can be unambiguously converted into Gurobi indices. In particular, \texttt{hard\_time\_constr} is injected as a hard constraint into the model, ensuring that no returned solution violates the availability configured in the frontend.

Finally, the backend response returns the computed schedule as a list of assignments containing \texttt{timeslot\_id}, \texttt{teacher\_id}, \texttt{course\_id}, and \texttt{room\_id}. Optionally, the canonical input including mapping information can be returned as well (\texttt{include\_input=true}), allowing the frontend to reconstruct grid positions and human-readable labels without additional database queries. Overall, this yields a compact yet robust integration between UI, persistence, and optimization despite the initial JS/Python interface friction.

\subsection{Results and Evaluation}
In summary, we have implemented software capable of generating schedules that satisfy all hard constraints when feasible, while also producing suboptimal solutions to accommodate soft constraints.

The software was tested on a limited dataset, in which it performed adequately. However, a significant increase in computational time was observed as problem size grew, and for datasets approaching realistic scale, we expect runtime to exceed 10 minutes or more.

In the absence of a complete dataset, weights for soft constraints were manually assigned, as there was no independent metric for schedule quality beyond the minimization objective. Ideally, with access to full data, it would be possible to calibrate these weights to best suit specific instances of the scheduling problem.

The current state of the project, and in particular its Gurobi-based component, leaves considerable room for future development. In particular, improvements could be made by incorporating user feedback, such as identifying missing but desired soft constraints and adjusting weights to better reflect practical priorities.
% ==========================================
% {Part III: Heuristic Optimization via Simulated Annealing
% ==========================================

\section{Part III: Heuristic Optimization via Simulated Annealing (SA)}
% Added this introduction paragrapgh. changed the hierarchy of the titles. The rest remain the same. Maybe you could also add a conclusion paragrapgh, and add some references from the papers we have read.

While the Gurobi solver guarantees the best possible schedule for smaller tasks, the Simulated Annealing (SA) module is a faster alternative for larger and more complex problems. SA doesn't guarantee a ``perfect'' mathematical solution, but it quickly finds high-quality schedules — much faster than trying every possible combination. This makes it perfect for the real-time updates in the OpTime frontend.

\subsection{Simulated Annealing Algorithm}

The physical process of Annealing in metallurgy inspires Simulated Annealing. When you heat metal and then cool it slowly, the atoms have enough energy to move around and settle into a low-energy, stable configuration. If you cool it too quickly, the atoms get stuck in a suboptimal arrangement.

In optimization, we use this analogy. The ``energy'' of our system is the cost of a schedule — i.e., the number of constraint violations it has. We start with a random schedule and make small changes to it. Here is the key insight: early in the process (at high temperature), we accept many changes, even those that make the schedule slightly worse. Starting with high temperature prevents us from getting stuck in a local optimum — a schedule that's better than all its neighbors but not the best possible schedule overall. As we continue (lowering temperature), we become more selective, mainly accepting improvements.

The temperature parameter $\tau(t)$ controls this behavior. At high temperatures, we explore the solution space widely. At low temperatures, we exploit reasonable solutions by making minor refinements.

\subsubsection{Four Required Components}

To apply Simulated Annealing to the timetabling problem, we must define four components: the solution space, the cost function, the neighbor relation, and the optimization goal.

\paragraph{Solution Space}

The solution space $C$ is the set of all possible complete assignments. A solution $S \in C$ is a complete assignment of all courses in $B$ to resources in $E$.

Let $B$ be the set of courses. Let $E = R \times T \times L$ where $R$ is the set of rooms, $T$ is the set of time slots, and $L$ is the set of lecturers.

A solution $S$ assigns to each course $b \in B$ a triple consisting of room, time slot, and lecturer. For example:
$$S(\text{Photogrammetric CV}) = (\text{SR\_A}, \text{Mon-09:15-10:45}, \text{Prof. Rodehorst})$$

This means Photogrammetric CV is scheduled in room SR\_A on Monday from 9:15 to 10:45, taught by Prof. Rodehorst.

\paragraph{Cost Function}

The cost function $c : C \to \mathbb{R}$ measures how bad a schedule is. Lower cost means better schedule. We define:
$$c(S) = w_1 \cdot f_1(S) + w_2 \cdot f_2(S) + ... + w_n \cdot f_n(S)$$

where $f_i$ are individual penalty functions and $w_i$ are weights representing constraint importance.

For our implementation:
$$c(S) = 100 \cdot (\text{room conflicts}) + 100 \cdot (\text{lecturer conflicts}) + 50 \cdot (\text{workload}) + 10 \cdot (\text{gaps}) + 10 \cdot (\text{Friday})$$

For example, if a schedule has 2 room conflicts, 1 Friday afternoon class, and 3 student gaps:
$$c(S) = 100 \times 2 + 10 \times 1 + 10 \times 3 = 240$$

The penalties are classified by severity:
\begin{itemize}
    \item Hard constraints (weight = 100): Must be satisfied. Examples: no room conflicts, no lecturer conflicts
    \item Medium constraints (weight = 50): Important quality indicators. Example: maximum daily workload
    \item Soft constraints (weight = 10): Preferences. Examples: minimize gaps, avoid Friday afternoons
\end{itemize}

\paragraph{Optimization Goal}

Our optimization problem is to find a solution $S^* \in C$ such that for all $S \in C$:
$$c(S^*) \leq c(S)$$

We call such $S^*$ an optimal solution.

\paragraph{Neighbor Relation}

We define a binary relation $\sim$ on $C \times C$. For $S, S' \in C$, we say $S \sim S'$ if and only if we can obtain $S'$ from $S$ by exactly one of the following operations.

\textbf{Technique 1: Construct and destroy}

We select one course $b \in B$. Let $e$ be the current resource assigned to $b$ in $S$. We choose a new resource $e' \in E$ and create:
$$S' = (S \setminus \{(b,e)\}) \cup \{(b,e')\}$$

We remove exactly one assignment (the old one for course $b$) and add exactly one assignment (the new one for course $b$).

Example: Move Photogrammetric CV from (SR\_A, Mon-09:15, Prof. Rodehorst) to (SR\_H, Tue-09:15, Prof. Rodehorst).

\textbf{Technique 2: Swap Operation}

We select two courses $b_1, b_2 \in B$ with $b_1 \neq b_2$, and their resources $e_1, e_2 \in E$ where $(b_1,e_1) \in S$ and $(b_2,e_2) \in S$, then:
$$S' = (S \setminus \{(b_1,e_1), (b_2,e_2)\}) \cup \{(b_1,e_2), (b_2,e_1)\}$$

Example: Swap Virtual Reality at (SR\_H, Wed-13:30) with HCI Theory at (LH\_HK7, Mon-13:30).

We define the neighborhood of solution $S$ as:
$$N(S) := \{S' \in C | S \sim S'\}$$

\subsubsection{The Algorithm}

The Simulated Annealing algorithm proceeds as follows:

\begin{verbatim}
Algorithm: MinimizationSimulatedAnnealing(I)
Input: instance I
Output: solution S

1: let S be an initial solution S_0
2: let k be a constant
3: let t = 0
4: while we continue do
5:     let S' in N(S) be chosen randomly
6:     if c(S') <= c(S) then
7:         S := S'
8:     else
9:         with probability e^(-(c(S')-c(S))/(k*tau(t))) let S := S'
10:        otherwise let S unchanged
11:    end if
12:    t := t + 1
13: end while
14: Return(S)
\end{verbatim}

\subsubsection{Acceptance Probability}

When we generate a neighbor schedule $S'$ from current schedule $S$, we calculate the cost difference:
$$\Delta = c(S') - c(S)$$

If $\Delta < 0$, the neighbor is better—we always accept it. If $\Delta > 0$, the neighbor is worse—we accept it with probability:
$$P(\text{accept}) = e^{-\Delta / (k \cdot \tau(t))}$$

where $k$ is a constant and $\tau(t)$ is the current temperature.

Example: Current schedule has cost 200, neighbor has cost 250, temperature is 100, and $k=1$. Then $\Delta = 50$ and $P(\text{accept}) = e^{-0.5} \approx 0.606$ (60\% chance). This probability decreases as temperature drops.

Our implementation guarantees that hard constraints are never violated by validating each neighbor during generation, ensuring only valid neighbors are considered.

\section{Conclusion}
OpTime started as a project for the Faculty of Media, but the system is general enough to be extended to other faculties or even the whole university with relatively small changes to the data model.

The workflow: Generate, Adjust manually, Improve keeps the process transparent and gives administrators control over the final result. The frontend hides the complexity of the two solvers and presents everything as a simple, visual interface.

% For future improvment, please check and feel free to add anything I might have missed.

For future work, the SA algorithm could be tuned further to get better results faster. The Gurobi model could also be tested on larger datasets beyond the Faculty of Media. On the frontend side, features like importing data directly from university systems would make the tool much more practical in a real administrative setting.

\printbibliography
\end{document}